\documentclass[10pt,leqno]{amsart}
\usepackage{graphicx}
\baselineskip=16pt

\usepackage{indentfirst,csquotes}

\topmargin= .5cm
\textheight= 20cm
\textwidth= 32cc
\baselineskip=16pt

\evensidemargin= .9cm
\oddsidemargin= .9cm

\usepackage{amssymb,amsthm,amsmath}
\usepackage{xcolor,paralist,hyperref,fancyhdr,etoolbox}
\usepackage[numbers,sort&compress]{natbib}
\newtheorem{theorem}{Theorem}[]
\newtheorem{definition}[theorem]{Definition}
\newtheorem{example}[theorem]{Example}
\newtheorem{lemma}[theorem]{Lemma}
\newtheorem{proposition}[theorem]{Proposition}
\newtheorem{corollary}[theorem]{Corollary}
\newtheorem{conjecture}[theorem]{Conjecture}

\hypersetup{ colorlinks=true, linkcolor=black, filecolor=black, urlcolor=black, citecolor=black }

\begin{document}

\title{Model Maju Dua-Zona \textit{Partial-LTE} dengan Transfer Radiatif
Deterministik untuk Rekonstruksi Spektrum Sintetis LIBS Bebas Kalibrasi}

\author[Walidain, Idris, Saddami]{Birrul Walidain, Nasrullah Idris, Khairun Saddami}
\date{\today}
\address{Department of Physics, Faculty of Mathematics and Natural Sciences,
Universitas Syiah Kuala, Banda Aceh, Indonesia}
\email{birrul.walidain@mhs.usk.ac.id}

\begin{abstract}
Analisis kuantitatif pada \textit{Calibration-Free Laser-Induced Breakdown Spectroscopy} (CF-LIBS) sangat bergantung pada asumsi plasma homogen berkesetimbangan termodinamika lokal (LTE). Namun, plasma transien memiliki gradien spasial dan opasitas optik yang ekstrem, memicu penyerapan mandiri (\textit{self-absorption}) yang mendistorsi spektrum emisi. Lebih jauh, model kinetika kolisional-radiatif (CR) murni sering gagal akibat ketiadaan data penampang kolisional yang lengkap pada basis data standar (fenomena ``Gurun Data Atomik''). Untuk mengatasi keterbatasan ini, penelitian ini mengusulkan model maju (\textit{forward model}) dua-zona (Core-Shell) deterministik yang menggabungkan \textit{partial}-LTE (pLTE) dengan Persamaan Transfer Radiatif (RTE). Paradigma pLTE menstabilkan populasi eksitasi internal dari defisit data kolisional, sementara RTE secara analitik menyelesaikan efek \textit{self-reversal} makroskopis. Dengan mengintegrasikan parameter pelebaran Stark secara eksak dari literatur dibandingkan menggunakan korelasi probabilistik, model ini secara akurat mereproduksi depresi spektral menyerupai garis Fraunhofer pada transisi resonansi yang tebal secara optik. Arsitektur deterministik ini sepenuhnya menghindari algoritma \textit{machine learning} kotak-hitam pada mesin fisika dasar, menyediakan generator spektrum sintetis bebas kalibrasi yang ketat secara fisis. Kerangka kerja ini mereduksi ketidakpastian kuantitatif secara signifikan dalam menganalisis material bermatriks kompleks, menawarkan landasan yang tangguh untuk diagnostik plasma dinamis.
\end{abstract}

\maketitle

\bigskip

\section{Posisi Riset, Kebaruan, dan Arah Pengembangan}

Manuskrip ini memosisikan model maju dua-zona \textit{partial}-LTE dengan transfer radiatif deterministik sebagai fondasi komputasi untuk tahap riset berikutnya, yaitu pengembangan \textit{Physics-Constrained Spectral AI} bagi analisis \textit{Laser-Induced Breakdown Spectroscopy} (LIBS). Arah pengembangan ini muncul dari kebutuhan praktis pada inspeksi material cepat, \textit{in-situ}, dan minim preparasi untuk sampel kompleks seperti tanah vulkanik, bijih, paduan logam, dan material industri strategis. Pada sistem semacam ini, ribuan garis emisi atomik dan ionik saling berhimpit, diperburuk oleh pelebaran Stark-Doppler dan \textit{self-absorption}, sehingga kuantifikasi unsur menjadi tidak stabil jika hanya bertumpu pada asumsi plasma homogen dan transparan optik.

Dalam lanskap metodologis LIBS, terdapat tiga keterbatasan utama yang memotivasi perluasan riset ini. Pertama, metode kurva kalibrasi membutuhkan \textit{matrix-matched standards} yang sering tidak tersedia pada kondisi lapangan. Kedua, \textit{Calibration-Free LIBS} klasik berbasis Saha--Boltzmann sangat sensitif terhadap validitas \textit{Local Thermodynamic Equilibrium} (LTE), sehingga mudah gagal pada plasma yang optis tebal. Ketiga, model \textit{deep learning} murni yang bersifat \textit{black box} cenderung hanya mempelajari korelasi statistik dangkal dan rentan kehilangan generalisasi saat distribusi data berubah. Oleh sebab itu, riset ini diarahkan untuk menjembatani keunggulan model fisika deterministik dengan efisiensi inferensi model pembelajaran mesin yang dikendalikan oleh kendala fisika.

Kebaruan konseptual penelitian ini terletak pada reposisi model maju deterministik bukan hanya sebagai alat simulasi, tetapi sebagai tulang punggung pembangkit data dan pembatas manifold fisis untuk arsitektur \textit{physics-informed learning}. Dalam peta jalan jangka menengah, fondasi ini diarahkan menuju \textit{Physics-Constrained Spectral Transformer} (PCST), yaitu arsitektur \textit{spectral transformer} yang dirancang untuk menangkap dependensi spektral jarak jauh pada kanal LIBS yang panjang, sambil tetap tunduk pada residu hukum transfer radiatif dan kendala opasitas. Dengan demikian, model fisika yang dibangun di manuskrip ini bukanlah artefak terpisah, melainkan blok dasar menuju sistem analisis LIBS yang cepat, kuantitatif, dan lebih dapat dipertanggungjawabkan secara fisis.

Secara \textit{state of the art}, pendekatan AI untuk LIBS dalam beberapa tahun terakhir didominasi oleh model \textit{data-driven} murni untuk klasifikasi atau regresi sederhana, sedangkan pendekatan berbasis fisika murni seperti pemecahan transfer radiatif komprehensif cenderung akurat tetapi mahal secara komputasi untuk aplikasi \textit{real-time}. Posisi riset ini berada di tengah dua kutub tersebut: model maju deterministik digunakan untuk menjaga kausalitas fisis, lalu hasilnya diproyeksikan ke arah arsitektur inversi cepat berbasis \textit{physics-informed surrogate} maupun \textit{spectral transformer}. Dengan framing ini, manuskrip ini menyumbang kebaruan pada level fondasi metodologis: ia menyediakan generator spektrum sintetis dua-zona yang siap dipakai sebagai basis pelatihan, validasi, dan pengembangan model AI terikat fisika pada tahap riset berikutnya.

Roadmap penelitian yang melandasi manuskrip ini dapat diringkas ke dalam empat fase. Fase pertama adalah era kemometrik klasik dan CF-LIBS berbasis asumsi homogen. Fase kedua adalah era \textit{deep learning} murni yang meningkatkan akurasi prediksi tetapi belum menyelesaikan persoalan kausalitas fisika dan \textit{self-absorption}. Fase ketiga, yang menjadi lingkup manuskrip ini, adalah pembangunan fondasi \textit{physics-informed machine learning} melalui model maju dua-zona, dataset sintetis, dan pipeline inversi awal yang masih terkendali secara fisis. Fase keempat adalah optimasi efisiensi, kompresi model, dan peluang hilirisasi menuju sistem LIBS cerdas yang lebih ringan untuk penggunaan lapangan.

% ============================================================================
\section{Introduction-fix}

Asumsi plasma homogen dalam Kesetimbangan Termodinamika Lokal (LTE) murni bertindak sebagai prasyarat fundamental bagi metode analisis kuantitatif \textit{calibration-free Laser-Induced Breakdown Spectroscopy} (CF-LIBS) \cite{favre_towards_2025, cristoforetti_local_2010}; namun, prasyarat teoretis ini terbukti gagal untuk mendeskripsikan dinamika rekombinasi plasma riil yang bergradien tinggi \cite{favre_merlin_2025, zaitsev_two-zone_2024}. Dalam kondisi ideal, kerangka termodinamika ini memastikan parameter fundamental plasma, yakni temperatur elektron ($T_e$) dan densitas elektron ($n_e$), dapat diekstraksi secara presisi guna mengkuantifikasi komposisi unsur secara langsung, yang juga menjadi landasan utama dalam pemodelan campuran ionik pada plasma astrofisika \cite{fritzsche_atomic_2025}. Sebaliknya, penyimpangan dari kondisi LTE ini tidak hanya mendistorsi akurasi analitik, melainkan mereduksi mekanisme kuantifikasi menjadi tidak valid \cite{cristoforetti_local_2010, hansen_modeling_2021, zaytsev_accurate_2026}. Mengingat kompleksitas tersebut, perumusan karakterisasi plasma yang tervalidasi secara ketat menjadi sangat esensial, khususnya pada analisis material bermatriks kompleks yang secara inheren rentan terhadap anomali ketidakhomogenan spasial saat pengamatan; seperti yang diaplikasikan pada instrumen eksplorasi antariksa \cite{sawyers_database_2025, manelski_libs_2024}, analisis batuan pada kondisi atmosfer ekstrem \cite{zhang_improving_2025}, karakterisasi \textit{in-situ} transisi fase logam cair \cite{leosson_analysis_2022}, evaluasi abrasi pelindung panas wahana saat memasuki atmosfer \cite{leistikow_quantification_2026}, maupun pada pemantauan \textit{in-situ} dinding reaktor fusi nuklir \cite{marin_roldan_libs_2021}.

Upaya konvensional untuk menjustifikasi ekuilibrium absolut tersebut secara tradisional bersandar pada kriteria McWhirter (1965); namun, perumusan klasik ini secara empiris terbukti menjadi syarat yang diperlukan tetapi tidak memadai untuk mendeskripsikan dinamika plasma riil \cite{bultel_mcwhirter_2025}. Kegagalan validasi ini berakar pada dua distorsi destruktif: secara spasial, gradien termal memicu efek penyerapan mandiri (\textit{self-absorption}) yang merusak linearitas profil spektral emisi \cite{tang_review_2024}; dan secara termodinamika, anomali populasi tingkat energi molekuler—seperti pada intensitas transisi vibrasi-rotasi—menunjukkan deviasi ekstrem dari prediksi LTE konvensional sehingga menuntut pendekatan multi-temperatur \cite{zaytsev_accurate_2026}. Untuk memitigasi anomali spasial ini, sejumlah studi komputasi mutakhir mulai mengintegrasikan arsitektur \textit{Artificial Neural Network} (ANN) dengan model dua-zona (\textit{two-zone model}) \cite{zaitsev_two-zone_2024, amogh_libs_2026,m_s_machine_2025}. Kendati pendekatan komputasi \textit{data-driven} ini berhasil mereduksi interferensi spasial, model-model tersebut mengalami kecacatan epistemologis karena masih tertahan pada asumsi kondisi stasioner. Padahal, kajian kinetika teoretis terbaru menegaskan bahwa plasma secara fundamental beroperasi secara sangat transien; di mana laju pendinginan yang masif dan tingginya opasitas radiatif ($\tau$) memaksa ambang batas densitas elektron melompat jauh melampaui prediksi stasioner \cite{bultel_mcwhirter_2025}. Kondisi ini mengekspos sebuah celah metodologis absolut: arsitektur komputasi yang semata-mata mengandalkan resolusi spasial tanpa intervensi penyelesaian kinetik transien berisiko kuat memberikan validasi keliru (\textit{false validation}) pada spektrum yang sejatinya non-LTE. 

Upaya komputasional termutakhir untuk merespons distorsi opasitas ini telah berevolusi, mulai dari penyelesaian analitik Persamaan Transfer Radiatif (RTE) \cite{favre_merlin_2025} hingga integrasi proksi \textit{machine learning} (ML) untuk mengakselerasi waktu komputasi \cite{favre_towards_2025}. Kendati inovasi metodologis ini berhasil mereduksi hambatan analitik secara signifikan, kerangka kerjanya secara umum masih bertumpu pada simplifikasi kondisi stasioner atau beroperasi sebagai ekstraktor emisi global tanpa resolusi spasial yang ketat. Keterbatasan dari asumsi stasioner ini semakin dipertegas oleh kajian teoretis terbaru yang menunjukkan bahwa perumusan ekuilibrium klasik—seperti kriteria McWhirter—cenderung mengunderestimasi ambang batas termodinamika pada plasma riil, terutama akibat tidak terakomodasinya laju pendinginan transien dan besarnya efek opasitas optik \cite{bultel_mcwhirter_2025}. Runtutan kompleksitas ini mengekspos sebuah celah metodologis yang esensial: pendekatan komputasi yang semata-mata mengandalkan pemecahan ruang spasial tanpa intervensi penyelesaian kinetik transien, maupun model prediktif yang mengabaikan kausalitas waktu, berisiko kuat memberikan validasi keliru terhadap spektrum plasma yang sejatinya berada pada rezim non-LTE.

Upaya konvensional untuk menjustifikasi ekuilibrium absolut tersebut secara tradisional bersandar pada kriteria McWhirter (1965); namun, perumusan klasik ini secara empiris terbukti menjadi syarat yang diperlukan tetapi tidak memadai untuk mendeskripsikan dinamika plasma riil \cite{bultel_mcwhirter_2025}. Kegagalan validasi ini berakar pada dua distorsi destruktif: secara spasial, gradien termal memicu efek penyerapan mandiri (\textit{self-absorption}) yang merusak linearitas profil spektral emisi \cite{tang_review_2024}; dan secara termodinamika, deviasi populasi tingkat energi molekuler menunjukkan anomali ekstrem dari prediksi konvensional sehingga menuntut pendekatan multi-temperatur pada kondisi non-LTE \cite{terashkevich_two-temperature_2024}. Untuk merespons anomali spasial dan distorsi opasitas ini, instrumentasi komputasional termutakhir telah berevolusi, mulai dari penyelesaian analitik Persamaan Transfer Radiatif (RTE) \cite{favre_merlin_2025}, hingga pengintegrasian proksi \textit{machine learning} (ML) dan \textit{Artificial Neural Network} (ANN) berbasis model dua-zona guna mengakselerasi waktu komputasi \cite{zaitsev_two-zone_2024, amogh_libs_2026, favre_towards_2025}. Kendati inovasi metodologis ini berhasil mereduksi interferensi spasial secara signifikan, kerangka kerjanya secara umum masih memiliki keterbatasan komputasional mendasar karena tertahan pada simplifikasi kondisi stasioner atau beroperasi sebagai ekstraktor emisi global tanpa resolusi spasial yang ketat. Keterbatasan asumsi stasioner ini bertabrakan dengan realitas kinetika teoretis terbaru yang menegaskan bahwa plasma secara fundamental bersifat sangat transien; di mana laju pendinginan yang masif dan tingginya opasitas radiatif ($\tau$) memaksa ambang batas termodinamika melompat jauh melampaui prediksi stasioner \cite{bultel_mcwhirter_2025}. Runtutan kompleksitas ini mengekspos sebuah celah metodologis absolut: arsitektur komputasi yang semata-mata mengandalkan resolusi spasial tanpa intervensi penyelesaian kinetik transien, maupun model prediktif yang mengabaikan kausalitas waktu, berisiko kuat memberikan validasi keliru (\textit{false validation}) pada spektrum yang sejatinya non-LTE.


Menghadapi celah metodologis tersebut sekaligus merespons trilema antara kegagalan ekuilibrium analitik, kelumpuhan komputasi deterministik, dan kesesatan kausalitas \textit{machine learning}, penelitian ini merumuskan sebuah resolusi arsitektur \textit{Physics-Informed Surrogate Model}. Arsitektur ini menolak secara fundamental penggunaan algoritma prediksi sebagai pencocok pola buta, melainkan memformulasikannya untuk beroperasi secara eksklusif di atas ruang termodinamika yang telah dikendalikan (\textit{pre-constrained manifold}). Untuk mencapai hal tersebut, fase komputasi dibelah menjadi dua ranah absolut: sebuah generator maju (\textit{forward model}) di fase \textit{offline} yang merajut termodinamika Dua-Zona \textit{partial}-LTE (pLTE) \cite{zaitsev_two-zone_2024} dengan penyelesaian analitik Persamaan Transfer Radiatif (RTE), serta sebuah proksi \textit{Support Vector Regression} (SVR) di fase \textit{online} \cite{wu_novel_2025}. Melalui pemisahan algoritmik yang deterministik ini, penelitian ini bertujuan untuk mengekstrak parameter suhu dan densitas plasma secara akurat dari spektrum emisi yang terdistorsi oleh \textit{self-absorption}, tanpa perlu menjalankan satu pun iterasi numerik RTE secara \textit{online}. Dalam integrasi hibrida ini, paradigma pLTE bertugas menstabilkan populasi eksitasi internal dari kegagalan ekuilibrium, sementara eksekusi SVR menjamin bahwa diagnostik tereksekusi secara seketika (\textit{real-time}) dengan tingkat kepatuhan absolut terhadap hukum fundamental mekanika kuantum; sebuah formulasi yang secara spesifik dirancang untuk menaklukkan anomali fisis pada batas optik tebal (\textit{optically thick limit}) sekaligus mengeliminasi hambatan waktu komputasional pada spektroskopi analitik.

Pada tingkat pengembangan lebih lanjut, arsitektur \textit{surrogate} ini tidak dibatasi pada SVR. Fondasi generator maju yang dibangun di manuskrip ini secara konseptual juga mendukung pengembangan model inversi yang lebih ekspresif, termasuk \textit{physics-informed neural network} dan \textit{spectral transformer}. Dalam kerangka tersebut, spektrum sintetis hasil model dua-zona dapat dipakai untuk membentuk pasangan data spektrum--parameter plasma, sementara residu transfer radiatif, opasitas, dan konsistensi rekonstruksi spektrum dapat dimasukkan langsung ke dalam fungsi rugi. Dengan kata lain, kontribusi manuskrip ini bersifat ganda: menyelesaikan persoalan forward modeling secara ketat, dan sekaligus menyediakan basis formal bagi arsitektur AI terikat fisika yang akan dikembangkan pada tahap riset lanjutan.


Ekspansi plasma yang ekstrem memicu gradien spasial dan pendinginan cepat, yang secara simultan menghancurkan ekuilibrium termodinamika tingkat dasar secara fundamental \cite{zaytsev_accurate_2026} dan memicu efek opasitas optik berupa penyerapan mandiri (\textit{self-absorption}) \cite{hansen_modeling_2021}. Pendekatan simplifikasi analitik yang buta terhadap transfer radiatif ini akan secara keliru menginterpretasikan pelebaran profil spektral akibat opasitas sebagai murni efek Stark kolisional, yang bermuara pada overestimasi densitas elektron ($n_e$) dan depresiasi ekstraksi suhu ($T_e$). Upaya untuk mengoreksi anomali opasitas tersebut menghadirkan kebuntuan metodologis baru. Pendekatan inversi deterministik murni—yang secara langsung mengeksekusi iterasi Persamaan Transfer Radiatif (RTE)—memang berhasil mempertahankan integritas termodinamika secara absolut \cite{favre_merlin_2025}. Namun, penggunaan algoritma pencarian akar non-linier untuk membalikkan matriks RTE memicu hambatan komputasi (computational bottleneck) yang ekstrem, sehingga menggugurkan kelayakan metode analitik ini untuk keperluan diagnostik analitik \cite{favre_towards_2025}. Di sisi ekstrem lainnya, substitusi komputasi menggunakan algoritma machine learning konvensional menawarkan kecepatan, tetapi menyimpan cacat epistemologis. Sebagaimana dievaluasi oleh Hao et al. \cite{hao_machine_2024}, mayoritas algoritma murni data-driven beroperasi sebagai 'kotak-hitam' (black-box) stokastik yang mencederai kausalitas mekanika kuantum, memutus rantai hukum termodinamika dasar, dan memicu ketergantungan irasional pada akuisisi dataset empiris berskala masif.

Menghadapi trilema antara kegagalan ekuilibrium analitik, kelumpuhan komputasi deterministik, dan kesesatan kausalitas \textit{machine learning}, penelitian ini merumuskan sebuah resolusi arsitektur \textit{Physics-Informed Surrogate Model}. Arsitektur ini menolak secara fundamental penggunaan \textit{machine learning} sebagai pencocok pola buta, melainkan memformulasikannya untuk beroperasi secara eksklusif di atas ruang termodinamika yang telah dikendalikan (\textit{pre-constrained manifold}). Untuk mencapai hal tersebut, fase komputasi dibelah menjadi dua ranah absolut: sebuah generator maju (\textit{forward model}) di fase \textit{offline} yang merajut termodinamika Dua-Zona \textit{partial}-LTE (pLTE) \cite{zaitsev_two-zone_2024} dengan penyelesaian analitik Persamaan Transfer Radiatif (RTE), serta sebuah proksi \textit{Support Vector Regression} (SVR) di fase \textit{online} \cite{wu_novel_2025}.

Melalui pemisahan algoritmik yang deterministik ini, penelitian ini bertujuan untuk mengekstrak parameter suhu dan densitas plasma secara akurat dari spektrum emisi yang terdistorsi oleh \textit{self-absorption}, tanpa perlu menjalankan satu pun iterasi numerik RTE secara \textit{online}. Dalam arsitektur hibrida ini, paradigma pLTE bertugas menstabilkan populasi eksitasi internal dari kegagalan ekuilibrium, sementara penggunaan SVR menjamin bahwa ekstraksi diagnostik tereksekusi secara seketika (\textit{real-time}) dengan tingkat kepatuhan absolut terhadap hukum termodinamika fundamental. Integrasi ini secara spesifik dirancang untuk menaklukkan anomali fisis pada batas optik tebal (\textit{optically thick limit}), sekaligus mengeliminasi hambatan waktu komputasional yang selama ini melumpuhkan analisis spektroskopi analitik berbasis plasma.

Asumsi plasma homogen yang selama ini digunakan dalam metode \textit{Calibration-Free Laser-Induced Breakdown Spectroscopy} (CF-LIBS) terbukti tidak memadai untuk mendeskripsikan dinamika plasma riil yang bergradien tinggi \cite{zaitsev_two-zone_2024}. Keandalan metode kuantitatif CF-LIBS mutlak bergantung pada pemenuhan syarat \textit{Local Thermal Equilibrium} (LTE) sebagai prasyarat utama \cite{cristoforetti_local_2010,favre_merlin_2025}. Dalam kondisi ideal, kerangka termodinamika ini memastikan parameter fundamental plasma, yakni temperatur elektron ($T_e$) dan densitas elektron ($n_e$), dapat diekstraksi secara presisi guna mengkuantifikasi komposisi unsur secara langsung, sebuah mekanisme fisis universal yang juga menjadi landasan utama dalam pemodelan campuran ionik pada plasma astrofisika \cite{fritzsche_atomic_2025}. Sebaliknya, penyimpangan dari kondisi LTE ini tidak hanya mendistorsi akurasi analitik, melainkan mereduksi mekanisme kuantifikasi menjadi tidak valid \cite{cristoforetti_local_2010,hansen_modeling_2021}. Mengingat kompleksitas pada karakterisasi plasma yang tervalidasi secara ketat menjadi sangat esensial, khususnya dalam analisis material matriks kompleks yang melekat rentan terhadap anomali ketidakhomogenan plasma pada pengamatan, seperti yang ditemui pada instrumen eksplorasi antariksa \cite{sawyers_database_2025,manelski_libs_2024}, analisis batuan pada kondisi atmosfer ekstrem \cite{zhang_improving_2025}, karakterisasi \textit{in-situ} transisi fase logam cair \cite{leosson_analysis_2022}, evaluasi abrasi pelindung panas wahana saat memasuki atmosfer \cite{leistikow_quantification_2026}, maupun pada pemantauan \textit{in-situ} dinding reaktor fusi nuklir \cite{marin_roldan_libs_2021}.

Meskipun kriteria McWhirter (1965) secara tradisional digunakan sebagai landasan validasi LTE, perumusan klasik ini terbukti menjadi syarat yang diperlukan namun tidak memadai untuk mendeskripsikan dinamika plasma riil \cite{bultel_mcwhirter_2025}. Distorsi paling destruktif pada plasma adalah dalam bentuk penyerapan mandiri (\textit{self-absorption}), di mana struktur spasial yang berlapis merusak linearitas profil spektral emisi \cite{tang_review_2024}. Ketidakpastian ini semakin diperumit oleh anomali pada profil populasi tingkat energi internal molekuler di mana konfigurasi intensitas transisi vibrasi-rotasi menunjukkan deviasi signifikan dari prediksi LTE tradisional dan menuntut pendekatan diagnostik dua-temperatur pada kondisi non-LTE \cite{terashkevich_two-temperature_2024}. Untuk memitigasi anomali spasial tersebut, studi-studi komputasi terkini beralih pada arsitektur \textit{Artificial Neural Network} (ANN) yang mengintegrasikan model dua zona (\textit{two-zone model}) \cite{zaitsev_two-zone_2024,amogh_libs_2026,m_s_machine_2025}. Namun, sementara pendekatan komputasi \textit{data-driven} ini berhasil mereduksi interferensi spasial, model-model tersebut masih tertahan pada penyederhanaan kondisi stasioner. Padahal, kajian kinetika teoretis terbaru membuktikan bahwa plasma secara fundamental bersifat sangat transien; laju pendinginan yang ekstrem dan tingginya opasitas radiatif ($\tau$) memaksa ambang batas densitas elektron minimal melompat jauh melampaui prediksi stasioner konvensional \cite{bultel_mcwhirter_2025}. Kondisi ini menyisakan sebuah celah metodologis, di mana arsitektur analitik yang hanya mengandalkan pemecahan ruang spasial tanpa melibatkan filtrasi kinetik berbasis waktu berisiko kuat memvalidasi spektrum non-LTE secara keliru.

Keterbatasan metode diagnostik konvensional seperti plot Saha-Boltzmann dalam menghadapi plasma transien telah diekspos secara empiris. Hansen \textit{et al.} \cite{hansen_modeling_2021} membuktikan bahwa asumsi plasma homogen gagal mengkompensasi distorsi opasitas, dan menegaskan bahwa penyederhanaan gradien spasial melalui arsitektur Transfer Radiatif (RTE) Dua-Zona adalah kerangka yang paling fisis dan akurat untuk memodelkan serapan mandiri (\textit{self-absorption}) tanpa batas hidrodinamika 3D. Lebih jauh, kegagalan ekuilibrium absolut di dalam plasma yang terekspansi dengan cepat juga telah dikonfirmasi oleh Zaitsev dan Terashkevich \cite{zaitsev_two-zone_2024}, yang menuntut terpisahnya derajat kebebasan fisis dan penerapan koreksi Kinetika Termodinamika (\textit{partial}-LTE). Oleh karena itu, penelitian ini merakit kedua fondasi fisik tersebut---RTE Dua-Zona dari Hansen sebagai penyelesai opasitas spasial dan Kinetika pLTE dari Zaitsev sebagai stabilisator eksitasi temporal---menjadi satu mesin generator deterministik. Pemisahan arsitektur ini melegitimasi penggunaan proksi interpolasi instan (\textit{Support Vector Regression}) di tahapan \textit{online}, selaras dengan pembuktian Liu \textit{et al.} \cite{lin_enhancing_2026} bahwa formulasi pencarian akar (\textit{root-finding}) berbasis komputasi heuristik atau optimasi multi-batasan ortodoks akan selalu lumpuh dan terjebak pada jebakan optimum lokal ketika dihadapkan pada redaman matriks inversi spektral.

Upaya ekstraksi suhu ($T_e$) dan densitas elektron ($n_e$) dari plasma transien yang terdistorsi oleh opasitas optik menghadirkan sebuah trilema metodologis yang melumpuhkan paradigma konvensional. Ketika dihadapkan pada spektrum cacat akibat penyerapan mandiri (\textit{self-absorption}), peneliti umumnya dipaksa memilih satu dari tiga jalur komputasi yang tersedia, yang masing-masing menyimpan cacat fisis atau arsitektural yang fatal.

Jalur pertama, yakni simplifikasi analitik yang diusung oleh model CF-LIBS tradisional, secara brutal melanggar realitas termodinamika. Pendekatan ini mengasumsikan transparansi optik absolut dan Kesetimbangan Termodinamika Lokal (LTE) murni \cite{cristoforetti_local_2010}. Akibatnya, pelebaran spektrum yang terjadi karena redaman opasitas secara keliru diinterpretasikan sebagai murni efek Stark kolisional, sehingga menghasilkan overestimasi densitas yang parah serta mengekstrak suhu dari tingkat energi yang sejatinya telah kehilangan ekuilibrium.

Jalur kedua, yakni inversi deterministik murni, mempertahankan integritas termodinamika namun berujung pada kelumpuhan komputasional. Menggunakan algoritma pencarian akar non-linear secara iteratif untuk membalikkan Persamaan Transfer Radiatif (RTE) memicu hambatan komputasi (\textit{computational bottleneck}) yang ekstrem \cite{rodriguez_mixkiprapcal_2021,espinosa_analysis_2019}. Menjalankan integrasi matriks yang kaku pada data eksperimen beresolusi tinggi memakan waktu yang sangat tidak rasional, sehingga seketika menggugurkan potensi diagnostik secara \textit{real-time}.

Jalur ketiga, yakni penerapan \textit{machine learning} konvensional (murni statistik), menawarkan kecepatan tinggi namun secara mutlak mencederai kausalitas mekanika kuantum \cite{m_s_machine_2025}. Algoritma ini terjebak dalam kesesatan kotak-hitam (\textit{black-box fallacy}), beroperasi murni sebagai pencocok pola tanpa pemahaman terhadap hukum termodinamika (Saha-Boltzmann). Ketergantungan absolut pada data eksperimen empiris membuat model ini rentan terhadap \textit{overfitting} dan melumpuhkan kemampuannya untuk melakukan generalisasi pada matriks material yang berbeda \cite{favre_towards_2025}.

Pendekatan simplifikasi analitik pada model CF-LIBS konvensional sering kali mendistorsi realitas termodinamika dengan memaksakan parameter transparansi optik absolut dan Kesetimbangan Termodinamika Lokal (LTE) murni secara homogen \cite{cristoforetti_local_2010}. Ketika asumsi idealis ini dipertahankan pada plasma transien, kalkulasi analitik murni terbukti memicu deviasi kuantitatif yang parah terhadap realitas fisis \cite{hansen_modeling_2021}. Secara mekanistik, redaman foton akibat opasitas optik (self-absorption) menyebabkan profil spektral mengalami pelebaran artifisial sekaligus penurunan intensitas puncak. Akibatnya, sistem diagnostik yang buta terhadap transfer radiatif akan menyalahartikan pelebaran ini sebagai murni efek Stark kolisional, yang bermuara pada overestimasi densitas elektron $n_e$. Secara bersamaan, depresiasi intensitas emisi tersebut memaksa algoritma mengekstrak suhu dari tingkat energi yang sejatinya telah kehilangan ekuilibrium, sehingga menghasilkan underestimasi suhu analitik yang keliru. 

Pendekatan inversi deterministik murni—yang secara langsung mengeksekusi iterasi Persamaan Transfer Radiatif (RTE)—memang berhasil mempertahankan integritas termodinamika secara absolut dalam memodelkan opasitas plasma \cite{hansen_modeling_2021, favre_merlin_2025}. Namun, paradigma ini secara inheren berujung pada kelumpuhan komputasional ketika diaplikasikan untuk diagnostik analitik. Sebagaimana yang dievaluasi oleh Favre et al. (2025), meskipun model komputasi radiatif murni sangat presisi, penggunaan algoritma pencarian akar non-linear untuk membalikkan matriks RTE secara iteratif memicu hambatan komputasi (computational bottleneck) yang ekstrem. Menjalankan beban integrasi numerik yang kaku pada spektrum eksperimen beresolusi tinggi menuntut biaya waktu yang tidak rasional, sebuah limitasi fundamental yang seketika menggugurkan kelayakan metode analitik ini untuk keperluan diagnostik secara seketika (real-time) \cite{favre_towards_2025}.

Meskipun substitusi analitik menggunakan algoritma \textit{machine learning} menjanjikan akselerasi komputasi yang signifikan, paradigma statistik murni ini menyimpan cacat epistemologis yang fatal jika diimplementasikan tanpa intervensi fisis. Sebagaimana diekspos secara komprehensif dalam tinjauan Hao et al. \cite{hao_machine_2024}, mayoritas algoritma \textit{data-driven} beroperasi sebagai kotak-hitam (\textit{black-box}) stokastik yang sepenuhnya mencederai kausalitas mekanika kuantum dan memutus rantai hukum termodinamika dasar. Keterputusan antara prinsip fisis dan arsitektur algoritma ini tidak hanya melumpuhkan kemampuan generalisasi model saat berhadapan dengan anomali efek matriks, tetapi juga memicu ketergantungan absolut pada akuisisi dataset empiris berskala masif yang secara eksperimental sangat tidak rasional. Berangkat dari limitasi absolut tersebut, penelitian ini menolak penggunaan \textit{machine learning} sebagai pencocok pola buta, dan memformulasikannya murni sebagai \textit{Surrogate Model} yang ruang pembelajarannya telah dibatasi dan didikte secara deterministik oleh generator fisika pLTE-RTE di fase \textit{offline}.

Menghadapi kebuntuan tiga arah ini, tujuan utama dari penelitian ini bukanlah merancang algoritma inversi LIBS konvensional, melainkan mengeksekusi resolusi tunggal yang mutlak terhadap Trilema Metodologis tersebut. Perancangan arsitektur \textit{Physics-Informed Surrogate Model} tidak lagi berstatus sebagai opsi analitik, melainkan satu-satunya entitas komputasi yang selamat. Integritas hukum mekanika kuantum sepenuhnya diamankan melalui generator \textit{partial}-LTE (pLTE) dan RTE \textit{offline} yang memproduksi anomali spektral deterministik melalui model dua-zona komprehensif \cite{zaitsev_two-zone_2024}, sementara hambatan komputasi (\textit{computational bottleneck}) dielakkan secara elegan dengan merekam kausalitas tersebut ke dalam proksi \textit{Support Vector Regression} (SVR) \textit{online}. Metodologi hibrida ini memastikan diagnostik ekstraksi termodinamika tereksekusi secara instan, kebal terhadap kesesatan \textit{black-box fallacy}, dan yang terpenting, tidak pernah mendistorsi apalagi mengkhianati fondasi fisika fundamentalnya \cite{favre_merlin_2025}.

% % ============================================================================


\section{Metode Komputasi: Arsitektur \textit{Physics-Informed Surrogate Model}}
% ============================================================================

Desain metodologis dalam penelitian ini dikonstruksi secara eksklusif menggunakan Arsitektur Hibrida Fase-Ganda (\textit{dual-phase hybrid architecture}). Pendekatan radikal ini secara ketat memisahkan proses pembangkitan data fisis (\textit{offline forward modeling}) di mesin fisika dasar (Subbab \ref{sec:rte} -- \ref{sec:stark}) dari proses ekstraksi parameter inversi seketika (\textit{online surrogate modeling}) pada mesin statistik intermediet (Subbab \ref{sec:surrogate}). Pemisahan arsitektural ini diformulasikan sebagai taktik komputasional esensial untuk memintas hambatan resolusi pada matriks Persamaan Transfer Radiatif (RTE), sekaligus mencabut sepenuhnya sifat stokastik \textit{black-box} dari algoritma prediksi konvensional.

\subsection{Tahapan Penelitian dan Integrasi dengan Pipeline Data}

Secara operasional, metode penelitian ini dibagi ke dalam tahapan berurutan yang menghubungkan model fisika, dataset sintetis, model inversi, dan validasi eksperimen. Tahap pertama adalah pembangunan model maju dua-zona \textit{partial}-LTE dan transfer radiatif, yang menjadi mesin simulasi utama untuk menghasilkan spektrum sintetik dengan efek \textit{self-absorption}. Tahap kedua adalah pembangkitan dataset sintetis skala besar pada rentang panjang gelombang 200--900 nm dengan konfigurasi komposisi dan parameter plasma yang bervariasi. Tahap ketiga adalah pelatihan model inversi berbasis \textit{physics-informed learning} pada spektrum hasil simulasi, baik tanpa seleksi fitur maupun dengan seleksi fitur \textit{minimum Redundancy Maximum Relevance} berbasis \textit{Mutual Information Quotient} (mRMR-MIQ). Tahap keempat adalah validasi pada sampel eksperimen tanah vulkanik yang memiliki referensi \textit{X-ray Fluorescence} (XRF) dan \textit{Calibration-Free LIBS} (CF-LIBS), sehingga model dapat dievaluasi terhadap acuan geokimia eksternal sekaligus baseline analitik LIBS klasik.

Dalam implementasinya, penelitian ini juga menggunakan pengelompokan dataset berdasarkan dua sumbu eksperimen. Sumbu pertama adalah kelompok komposisi: kelompok A untuk unsur mayor dan kelompok B untuk unsur mayor ditambah unsur jejak. Sumbu kedua adalah tier resolusi pragmatis: rendah, menengah, dan tinggi. Kombinasi keduanya membentuk dataset kode A-L, A-M, A-H, B-L, B-M, dan B-H. Desain ini memungkinkan evaluasi terkontrol terhadap pengaruh kompleksitas komposisi, kepadatan kanal spektral, dan strategi seleksi fitur terhadap performa model inversi.

Luaran dari tahapan ini tidak hanya berupa model terlatih, tetapi juga produk antara yang bernilai ilmiah tinggi: generator spektrum sintetis, dataset HDF5 berlabel termodinamika dan komposisi, model inversi \textit{physics-informed}, serta laporan evaluasi kuantitatif yang dapat langsung dibandingkan dengan data eksperimen. Dengan demikian, metode yang dibangun tetap sejalan dengan kebutuhan validasi laboratorium dan target pengembangan perangkat lunak riset pada tingkat kesiapan teknologi menengah.

\subsection{Arsitektur Geometri Dua-Zona dan Integrasi RTE} \label{sec:rte}

Untuk menangkap heterogenitas spasial dari fenomena penyerapan mandiri tanpa memicu biaya komputasi eksponensial dari simulasi hidrodinamika 3D penuh, plasma didiskretisasi secara spasial menjadi dua wilayah berbeda: inti emisif (\textit{Core}, $V_1$) dan selubung absorptif (\textit{Shell}, $V_p$) \cite{zaitsev_two-zone_2024,hansen_modeling_2021}. Transportasi foton melalui medium bergradien ini diatur oleh persamaan diferensial fundamental:
\begin{equation} \label{eq:rte_general}
\frac{dI_\lambda}{ds} = j_\lambda - \kappa_\lambda I_\lambda
\end{equation}

Untuk skema Dua-Zona, diasumsikan zona inti memancarkan radiasi primer $I_{\text{core}}(\lambda)$. Radiasi ini melintasi zona selubung yang memiliki ketebalan geometris $d_{\text{shell}}$ dan parameter optik konstan lokal ($j_{\text{shell}}$ dan $\kappa_{\text{shell}}$). Dengan mengintegrasikan kedua sisi melintasi selubung dari $s = 0$ hingga $s = d_{\text{shell}}$:
\begin{equation} \label{eq:rte_integral}
\int_{I_{\text{core}}}^{I_{\text{obs}}} \frac{dI_\lambda}{j_{\text{shell}} - \kappa_{\text{shell}} I_\lambda} = \int_{0}^{d_{\text{shell}}} ds
\end{equation}

Penyelesaian integral analitik tersebut menghasilkan rumusan mutlak yang mensintesis spektrum observasi final tanpa memerlukan asimilasi fungsi Benda Hitam Planck \cite{espinosa_analysis_2019}:
\begin{equation} \label{eq:rte_solution}
I_{\text{obs}}(\lambda) = I_{\text{core}}(\lambda) e^{-\tau_\lambda} + \frac{j_{\text{shell}}(\lambda)}{\kappa_{\text{shell}}(\lambda)} \left( 1 - e^{-\tau_\lambda} \right)
\end{equation}
Di mana kedalaman optik dievaluasi secara deterministik sebagai $\tau_\lambda = \kappa_{\text{shell}} d_{\text{shell}}$, dan suku $\frac{j}{k}$ merupakan rasio termodinamika kolisional-radiatif yang menolak berlakunya Hukum Kirchhoff pada rezim transien \cite{rodriguez_mixkiprapcal_2021}.

\subsection{Kinetika Populasi Deterministik: Paradigma \textit{Partial}-LTE}

Divergensi fundamental dari metodologi ini terhadap model homogen konvensional terletak pada perhitungan populasi tingkat energi ($n_i$). Dalam konteks plasma LIBS, evolusi temporal dari populasi pada tingkat energi ke-$i$ secara ideal diatur oleh sistem persamaan diferensial berpasangan (\textit{coupled rate equations}) Kolisional-Radiatif (CR):
\begin{equation} \label{eq:master_kinetic}
\frac{dn_i}{dt} = \sum_{j > i} \left( n_j A_{ji} + n_j n_e q_{ji}^{\text{down}} \right) + \sum_{j < i} n_j n_e q_{ji}^{\text{up}} - n_i \sum_{j < i} \left( A_{ij} + n_e q_{ij}^{\text{down}} \right) - n_i \sum_{j > i} n_e q_{ij}^{\text{up}}
\end{equation}
Di mana $A_{ji}$ adalah probabilitas emisi spontan Einstein, $n_e$ adalah densitas elektron, serta $q^{\text{up}}$ dan $q^{\text{down}}$ adalah koefisien laju eksitasi dan de-eksitasi kolisional.

Namun, efektivitas model CR murni secara fundamental dibatasi oleh fenomena ``Gurun Data Atomik'' (\textit{Atomic Data Desert}), di mana ketersediaan data penampang kolisional dalam basis data NIST tidak mencukupi untuk mendukung perhitungan CR yang deterministik pada seluruh tingkat energi \cite{bultel_mcwhirter_2025,favre_merlin_2025}. Akibatnya, matriks laju $\mathbf{M}$ dalam formulasi $\frac{d\mathbf{n}}{dt} = \mathbf{M} \cdot \mathbf{n}$ menjadi sangat jarang (\textit{sparse}), menyebabkan tingkat energi intermediet kelaparan data dan menghasilkan distribusi populasi yang secara numerik tidak stabil.

Untuk mengatasi degradasi kinetika ini, penelitian ini menerapkan paradigma \textit{Partial}-LTE (pLTE). Dalam paradigma pLTE, keseimbangan ionisasi (\textit{bound-free transitions}) diasumsikan telah mencapai termalisasi akibat frekuensi tumbukan elektron yang sangat ekstrem pada fase transien awal ekspansi plasma \cite{cristoforetti_local_2010}. Oleh karena itu, sebelum populasi internal dievaluasi, distribusi populasi makroskopis antara atom netral ($n_I$) dan ion terionisasi tunggal ($n_{II}$) pada setiap zona spasial dikalkulasi secara deterministik menggunakan Persamaan Saha:
\begin{equation} \label{eq:saha}
\frac{n_{e} n_{II}}{n_{I}} = 2 \frac{Z_{II}(T_e)}{Z_{I}(T_e)} \left( \frac{2\pi m_e k_B T_e}{h^2} \right)^{3/2} \exp\left(-\frac{E_{ion}}{k_B T_e}\right)
\end{equation}
di mana $Z(T_e)$ merupakan fungsi partisi bergantung suhu, $h$ adalah konstanta Planck, dan $E_{ion}$ adalah energi ionisasi elemen terkait \cite{bultel_mcwhirter_2025}. Setelah fraksi spesies makroskopis ini diekstraksi secara presisi dari Persamaan \ref{eq:saha}, populasi internal untuk setiap tingkat eksitasi terikat (\textit{bound excited states}) didistribusikan secara statistik menggunakan probabilitas Boltzmann untuk mengkompensasi ketiadaan data penampang kolisional:
\begin{equation} \label{eq:boltzmann_plte}
n_i = \frac{n_{\text{total}} \cdot g_i \exp\left(-E_i / k_B T_e\right)}{Z(T_e)}
\end{equation}
Di mana $g_i$ adalah berat statistik (\textit{degeneracy}), $E_i$ adalah energi tingkat, $k_B$ adalah konstanta Boltzmann, $T_e$ adalah suhu elektron, dan $Z(T_e) = \sum_i g_i \exp(-E_i/k_B T_e)$ adalah fungsi partisi. Pendekatan ini valid karena pada densitas plasma LIBS yang tinggi ($n_e \ge 10^{16}$ cm$^{-3}$), laju tumbukan elektron cukup cepat untuk metermalisasi populasi internal menuju kesetimbangan lokal, meskipun keseimbangan ionisasi antar spesies (yang memerlukan Persamaan Saha) belum tentu tercapai \cite{cristoforetti_local_2010,bultel_mcwhirter_2025}. Dengan demikian, paradigma pLTE secara eksplisit mempertahankan fisika transfer radiatif dua-zona untuk efek makroskopis (\textit{self-absorption}), sambil menstabilkan perhitungan populasi mikroskopis dari ketidaklengkapan basis data atomik.

\subsection{Koefisien Optik: Suku Emisi dan Absorpsi}

Koefisien emisi ($j_\lambda$) dan koefisien absorpsi ($\kappa_\lambda$) mendefinisikan secara kuantum seberapa kuat plasma memancarkan dan menyerap cahaya pada setiap panjang gelombang. Transisi ini didorong oleh emisi de-eksitasi dari tingkat atas ($n_k$) dan absorpsi dari tingkat bawah ($n_i$) \cite{favre_merlin_2025}:
\begin{equation} \label{eq:emission_coeff}
j_\lambda = \sum_{k \to i} \frac{hc}{4\pi\lambda} A_{ki} n_k \phi_V(\lambda)
\end{equation}

Dalam plasma LIBS berdensitas tinggi pada fase transien awal, foton yang dipancarkan dapat memicu elektron lain untuk turun tingkat (emisi terstimulasi). Oleh karena itu, bentuk final dari koefisien absorpsi mencakup koreksi keseimbangan detail (\textit{detailed balance}) \cite{espinosa_analysis_2019}:
\begin{equation} \label{eq:absorption_coeff}
\kappa_\lambda = \sum_{i \to k} \frac{\lambda^4}{8\pi c} \left(\frac{g_k}{g_i}\right) A_{ki} n_i \phi_V(\lambda) \left[ 1 - \frac{n_k g_i}{n_i g_k} \right]
\end{equation}
Faktor dalam kurung siku merekam secara absolut penyimpangan sistem dari ekuilibrium (\textit{Non-LTE departure}). Profil garis $\phi_V(\lambda)$ merupakan konvolusi Voigt yang dibahas pada subbab berikutnya.

\subsection{Kausalitas Pelebaran Garis Spektral}

Profil garis emisi dan absorpsi dalam plasma LIBS tidak bersifat monokromatik singular, melainkan mengalami pelebaran fisis yang direpresentasikan oleh fungsi konvolusi Voigt $\phi_V(\lambda)$ \cite{tang_review_2024}. Pelebaran ini merupakan akumulasi dari berbagai interaksi mikroskopis dan makroskopis di dalam plasma, yang secara taksonomi termodinamika dibagi menjadi pelebaran berprofil Gaussian dan Lorentzian.

\subsubsection{Pelebaran Gaussian (Doppler)}
Gerak termal acak dari partikel emitor relatif terhadap pengamat (spektrometer) menghasilkan pergeseran frekuensi relativistik berorde pertama. Berdasarkan distribusi kecepatan Maxwell-Boltzmann pada suhu ion/atom ($T_i$), profil yang dihasilkan adalah Gaussian murni dengan \textit{Full Width at Half Maximum} (FWHM) sebesar:
\begin{equation} \label{eq:broadening_doppler}
\Delta\lambda_D = \lambda_{ki} \sqrt{8 \ln 2 \frac{k_B T_i}{m_k c^2}}
\end{equation}
Di mana $\lambda_{ki}$ adalah panjang gelombang pusat transisi dan $m_k$ adalah massa partikel emitor.

\subsubsection{Pelebaran Lorentzian (Stark Kolisional)} \label{sec:stark}
Mekanisme pelebaran secara matematis menghasilkan profil Lorentzian, di mana besaran FWHM kumulatifnya umumnya didominasi secara mutlak oleh Pelebaran Stark Kolisional pada rezim LIBS ($n_e \ge 10^{16}$ cm$^{-3}$) \cite{tang_review_2024}.

Pelebaran Stark Kolisional disebabkan oleh gangguan medan listrik mikro dari partikel bermuatan (elektron dan ion) di sekitar emitor. FWHM Stark secara umum diformulasikan sebagai:
\begin{equation} \label{eq:broadening_stark}
\Delta\lambda_S = 2 \omega_{ki}(T) \frac{n_e}{n_{e,r}} \left[ 1 + 1.75 \beta(n_e, T) \right]
\end{equation}
Di mana $n_{e,r}$ adalah densitas elektron referensi, $\omega_{ki}(T)$ adalah parameter pelebaran tumbukan elektron, dan $\beta(n_e, T)$ adalah faktor kontribusi interaksi ion statis. Untuk plasma transien dengan densitas elektron tinggi ($n_e > 10^{16}$ cm$^{-3}$), suku kontribusi ion statis dapat diabaikan, sehingga persamaan menyusut menjadi dependensi linear murni terhadap densitas elektron ($n_e$):
\begin{equation} \label{eq:broadening_stark_simp}
\Delta\lambda_S \approx 2 \omega_{ki}(T) \frac{n_e}{n_{e,r}}
\end{equation}

Dalam arsitektur analitik ini, nilai parameter pelebaran $\omega_{ki}$ tidak diturunkan secara teoretis, melainkan diekstraksi secara eksak dari basis data empiris Griem (1974) \cite{griem_spectral_1974}.

Berbeda dengan metode komputasi konvensional analitik komprehensif (seperti kode MERLIN \cite{favre_merlin_2025}) yang sering memaksakan penggunaan aproksimasi korelasi polinomial saat parameter $\omega_{ki}$ absen---sebuah kompromi probabilistik yang rentan memicu \textit{underestimation} atau \textit{overestimation} profil spektral---penelitian ini secara fundamental menolak segala bentuk tebakan stokastik tersebut. Pemodelan fisis difokuskan secara eksklusif dan hanya dieksekusi pada ``Garis Diagnostik Utama'' (seperti doublet resonansi Ca II atau Si I) yang parameter mekanika kuantumnya telah ditranskripsi dan divalidasi secara absolut. Prinsip Ketidaktahuan Sokratik (\textit{Socratic Ignorance}) ini ditegakkan secara ekstrem demi menjaga ruang pengukuran sintetik dari polusi data fiktif, menjamin bahwa penampang densitas elektron ($n_e$) diekstraksi tanpa residu korelasi matematis buatan \cite{tang_review_2024, manelski_libs_2024}.

\subsubsection{Reduksi Deterministik pada Rezim LIBS Transien}
Berdasarkan analisis komparatif besaran pelebaran (\textit{broadening magnitude hierarchy}), arsitektur komputasi deterministik dalam penelitian ini melakukan pemangkasan fisis secara rasional. Pembentukan profil Voigt $\phi_V(\lambda)$ difokuskan murni pada konvolusi eksak antara pelebaran Doppler (Persamaan \ref{eq:broadening_doppler}) yang sensitif terhadap suhu, dan pelebaran Stark kolisional (Persamaan \ref{eq:broadening_stark_simp}) yang sensitif terhadap densitas:
\begin{equation} \label{eq:voigt_final}
\phi_V(\lambda) = \int_{-\infty}^{\infty} G(\lambda', \Delta\lambda_D) \cdot L(\lambda - \lambda', \Delta\lambda_S) d\lambda'
\end{equation}
Di mana hasil akhir spektrum emisi komputasional kemudian dikonvolusikan secara eksternal dengan Fungsi Aparatus (\textit{Instrumental Broadening}, $\Delta\lambda_A$) untuk mereplikasi batas resolusi optik dari perangkat keras spektrometer empiris \cite{favre_merlin_2025}.

\subsubsection{Fungsi Aparatus dan Redaman Distorsi Fraunhofer}
Spektrum yang dihasilkan dari integrasi RTE Dua-Zona (Persamaan \ref{eq:rte_solution}) secara teoretis memprediksi depresi Fraunhofer (\textit{self-reversal}) yang sangat ekstrem dan tajam pada garis-garis resonansi beropasitas tinggi. Namun, dalam observasi eksperimental murni, ketajaman profil fisis ini akan mengalami redaman (\textit{smearing effect}) secara signifikan akibat keterbatasan resolusi optik dari kisi difraksi dan celah (\textit{slit}) pada perangkat keras spektrometer, seperti sistem Echelle maupun \textit{Optical Multichannel Analyzer} (OMA) \cite{zaitsev_two-zone_2024}.

Untuk mereplikasi distorsi observasi ini secara deterministik, intensitas komputasional murni dari plasma ($I_{obs}$) wajib dikonvolusikan secara eksternal dengan Fungsi Aparatus berprofil Gaussian ($G_{inst}$):
\begin{equation} \label{eq:instrumental}
I_{final}(\lambda) = I_{obs}(\lambda) \otimes G_{inst}(\lambda, \Delta\lambda_{inst})
\end{equation}
di mana $\otimes$ menyatakan operator konvolusi diskrit, dan parameter resolusi aparatus ($\Delta\lambda_{inst}$) ditetapkan pada batas fisis perangkat spektrometer eksperimental (dalam penelitian ini, ditetapkan pada FWHM $\approx 0.1$ nm). Manuver konvolusi eksternal inilah yang secara fisis mendistribusikan ulang energi spektral dari sayap profil (\textit{line wings}) ke pusat cekungan, meratakan jurang \textit{self-absorption} ekstrem menjadi profil cekungan landai yang secara visual tervalidasi dan identik dengan spektrum pengamatan laboratorium \cite{favre_merlin_2025}.

\subsection{Fase Eksekusi Inversi Instan melalui Arsitektur \textit{Physics-Informed Surrogate Model}} \label{sec:surrogate}

Konklusi desain dari arsitektur metodologis ini bermuara pada persilangan dua dunia komputasi yang direalisasikan ke dalam dua fase sekuensial absolut:

\textbf{1. Deklarasi Fase \textit{Offline (Pre-Constrained Manifold)}.} Seluruh rangkaian formulasi analitik (dari Persamaan Saha, Boltzmann, hingga integrasi deterministik RTE dan pelebaran Voigt) dikompilasi secara rigid sebagai Generator fisis murni. Menolak menginterpolasi tingkat batas energinya \cite{favre_towards_2025}, generator mengeksekusi sapuan termodinamika masif ($T_e \in [4000, 20000]$\,K dan $n_e \in [10^{14}, 10^{18}]$\,cm$^{-3}$) untuk membombardir arsitektur algoritma statistik dengan ribuan spektrum. Keluaran dari fase ini bukanlah sekumpulan data buta, melainkan entitas data yang telah terkunci secara natif ke dalam ruang metrik yang mutlak mematuhi hukum termodinamika dan opasitas radiatif (\textit{pre-constrained physical manifold}).

\textbf{2. Deklarasi Fase \textit{Online (Instant Surrogate Proxies)}.} Pada tahap awal, arsitektur \textit{surrogate} dapat diwujudkan melalui model regresi kernel atau model saraf terikat fisika untuk memetakan spektrum ke parameter termodinamika dan komposisi. Dalam evolusi berikutnya, kerangka ini diarahkan ke \textit{Physics-Constrained Spectral Transformer} (PCST), yakni model transformer spektral yang dirancang untuk memproses korelasi kanal panjang dengan biaya komputasi yang masih rasional melalui mekanisme perhatian jarang (\textit{sparse attention}) seperti ProbSparse/Informer. Berbeda dari transformer murni yang hanya memaksimalkan kesesuaian statistik, PCST yang diusulkan akan menempatkan residu transfer radiatif, kendala opasitas, dan konsistensi rekonstruksi spektrum sebagai komponen fungsi rugi, sehingga model tidak sekadar mempelajari pola, tetapi juga mempertahankan keterikatan pada hukum fisika plasma.

Dalam konfigurasi ini, fase \textit{online} tidak lagi dipahami sebagai proses fitting iteratif yang mahal, melainkan sebagai inferensi cepat di atas manifold yang telah dibatasi oleh generator dua-zona. Oleh sebab itu, target utama bukan sekadar akurasi prediksi, tetapi pencapaian kompromi ilmiah antara ketepatan kuantifikasi, kepatuhan fisika, dan efisiensi inferensi. Arah ini sejalan dengan tujuan validasi riset pada sampel tanah vulkanik nyata yang memiliki referensi XRF dan CF-LIBS, serta membuka jalan bagi pengembangan perangkat lunak analisis LIBS yang dapat dipakai secara lebih praktis di laboratorium maupun lapangan.

% ============================================================================
\bibliographystyle{elsarticle-num}
\bibliography{reff}

\end{document}
